\begin{longtable}{p{3cm}p{10cm}}
AAP&Pädeva asutuse juurdepääsupunkt (\textit{Authority Access Point})\\
AP&Juurdepääsupunkt (\textit{Access Point})\\
API&Rakendusliides (\textit{Application Programming Interface})\\
AS4&Turvalise B2B sõnumivahetuse standardprofiil (\textit{Applicability Statement 4})\\
BDXL&Äridokumentide metaandmete asukohateenuse spetsifikatsioon (\textit{Business Document Metadata Service Location})\\
BDXR&Äridokumentide vahetusvõrgustike spetsifikatsioonide raamistik (\textit{Business Document Exchange Network})\\
CA&Sertifitseerimiskeskus (\textit{Certificate Authority})\\
CEF&Euroopa ühendamise rahastu (\textit{Connecting Europe Facility})\\
CNAME&DNS-i aliaskirje, mis suunab ühe nime teisele nimele (\textit{Canonical Name Record})\\
DNS&Domeeninimede süsteem (\textit{Domain Name System})\\
DNS-OARC&DNS-i operatsioonide, analüüsi ja teadusuuringute keskus (\textit{Domain Name System Operations, Analysis, and Research Center})\\
DNSSEC&DNS turbelaiendus (\textit{Domain Name System Security Extensions})\\
eFTI&Elektrooniline veoinfo (\textit{Electronic Freight Transport Information})\\
eFTI4EU&eFTI4EU projekt (\textit{Electronic Freight Transport Information for the European Union})\\
HTTP&Hüperteksti edastusprotokoll (\textit{Hypertext Transfer Protocol})\\
HTTPS&Turvaline hüperteksti edastusprotokoll (\textit{Hypertext Transfer Protocol Secure})\\
IP&Internetiprotokoll (\textit{Internet Protocol})\\
MD5&Krüptograafiline räsialgoritm MD5 (\textit{Message Digest Algorithm 5})\\
NAPTR&DNS-i kirjetüüp teenuse asukoha leidmiseks (\textit{Name Authority Pointer})\\
NIIS&Põhjamaade Interoperatiivsuse Instituut (\textit{Nordic Institute for Interoperability Solutions})\\
OASIS&Rahvusvaheline avatud IT-standardite arendusorganisatsioon (\textit{Organization for the Advancement of Structured Information Standards})\\
PEPPOL&Üleeuroopaline e-hangete andmevahetusvõrgustik (\textit{Pan-European Public Procurement OnLine})\\
POC&Kontseptsiooni tõestus (\textit{Proof of Concept})\\
REST&Ressursipõhine veebiarhitektuuri stiil (\textit{Representational State Transfer})\\
RFC&IETF tehniliste spetsifikatsioonide dokumendisari (\textit{Request for Comments})\\
SHA-256&SHA-2 perekonna 256-bitine räsialgoritm (\textit{Secure Hash Algorithm 256})\\
SML&Teenuse metaandmete leidmise komponent (\textit{Service Metadata Locator})\\
SMP&Teenuse metaandmete avaldamise komponent (\textit{Service Metadata Publisher})\\
SOAP&XML-põhine veebiteenuste sõnumivahetuse protokoll (\textit{Simple Object Access Protocol})\\
SQL&Struktureeritud päringukeel (\textit{Structured Query Language})\\
TLS&Transpordikihi turbeprotokoll (\textit{Transport Layer Security})\\
URL&Veebiressursi aadress (\textit{Uniform Resource Locator})\\
VoIP&IP-võrgu kaudu toimuv kõneside (\textit{Voice over Internet Protocol})\\
WS&Veebiteenus (\textit{Web Service})\\
WSDL&Veebiteenuse kirjelduse keel (\textit{Web Services Description Language})\\
XML&Laiendatav märgistuskeel (\textit{eXtensible Markup Language})\\
XMLDSIG&XML-vormingus digiallkirja standard (\textit{XML Digital Signature})\\
\end{longtable}
\addtocounter{table}{-1}
